\documentclass[11pt]{article}
\usepackage[margin=1in]{geometry}
\usepackage{amsmath,amssymb,amsthm}
\usepackage{graphicx}
\usepackage{placeins}
\usepackage{booktabs}
\usepackage[expansion=false]{microtype}
\usepackage[hidelinks]{hyperref}
\setlength{\emergencystretch}{3em}

\theoremstyle{definition}
\newtheorem{definition}{Definition}
\newtheorem{proposition}{Proposition}

\newcommand{\Sing}{\mathfrak{S}}

\title{\bf Prime Arithmetic Progressions as a Lattice Resonance:\\
the Admissibility Selection Rule and the\\
Singular-Series Enhancement Spectrum%
\thanks{Code and measured data: \url{https://github.com/Ruqing1963/6N-ap-resonance}}}
\author{Ruqing Chen\\
\small GUT Geoservice Inc., Montreal, Quebec, Canada\\
\small \texttt{ruqing@hotmail.com}}
\date{June 8, 2026}

\begin{document}
\maketitle

\begin{abstract}
We measure the count of $k$-term arithmetic progressions of primes as a function of the common
difference $\Delta$, on the $6N$ skeleton, sieving all primes to $X=10^{8}$. Viewed as a function of
$\Delta$ the count $N_k(\Delta)$ behaves exactly like a diffraction spectrum. (i) A sharp
\emph{selection rule}: $N_k(\Delta)$ is identically zero unless $\Delta$ is divisible by the primorial
$\prod_{p\le k}p$. We verify this as an absolute extinction---scanning every $\Delta\le1512$ for
$k=3$, the number of admissible-looking off-lattice progressions found is exactly zero. (ii) On the
admissible lattice $\Delta\in\mathrm{primorial}(k)\cdot\mathbb{Z}$, a clean \emph{resonance
enhancement}: each additional prime $q>k$ dividing $\Delta$ multiplies the count by the universal
factor $(q-1)/(q-k)$. This is the relative Hardy--Littlewood singular series
$\Sing_k(\Delta)=\prod_p\bigl(1-\nu_p(\Delta)/p\bigr)/(1-1/p)^{k}$, with
$\nu_p(\Delta)=1$ if $p\mid\Delta$ and $\min(k,p)$ otherwise; we verify it against the data to
$\lesssim1\%$ across $k=3,4,5$ (e.g.\ at $k=3$, $\Delta=2310$ gives a measured $3.768$ against the
predicted $3.750$; at $k=4$, $\Delta=210$ gives $8.11$ against $8.00$). Both the extinction rule and
the $(q-1)/(q-k)$ enhancement are exact consequences of admissibility; what the experiment adds is the
spectral presentation and the high-precision verification at $X=10^{8}$, together with the unified
per-prime resonance law across $k$. We make no claim about the existence of arbitrarily long
progressions: the Green--Tao theorem is the (untouchable) backdrop, not a target. As throughout, this
is a measurement, not a theorem; the Hardy--Littlewood heuristic is taken as input.
\end{abstract}

\section{Introduction}
That the primes contain arbitrarily long arithmetic progressions is the Green--Tao theorem
\cite{GT}, one of the central results of the century, proved by ergodic and transference methods whose
arithmetic content is famously hard to picture. The infinitude of three-term progressions is older,
due to van der Corput \cite{vdC}; the conjectural asymptotic count of $k$-term progressions is the
Hardy--Littlewood prime-tuple heuristic \cite{HL} specialised to the constellation
$\{0,\Delta,2\Delta,\dots,(k-1)\Delta\}$.

This programme has spent its recent papers measuring the local structure of one constellation after
another: the twin pair (Parts I, XXIII), the $\omega$-conditioned second moment (Part XXVI), the
twin-enrichment closed form (Part XXVIII), and the twin pair correlation governing the short-interval
variance (Part XXIX) \cite{PXXIX}. The pair correlation of Part XXIX is the $4$-tuple singular series
read as a function of separation. The present paper takes the natural next step: from two points to a
$k$-point constellation strung along an arithmetic progression, read as a function of the common
difference $\Delta$.

The result is best stated in the language of diffraction. Hold $k$ fixed and vary $\Delta$. The count
$N_k(\Delta)$ of $k$-term prime progressions with common difference $\Delta$ and first term up to $X$
shows two features of a lattice spectrum: a hard \emph{extinction rule} that forbids all $\Delta$ off
a primorial sublattice, and a \emph{resonance enhancement} on that sublattice with sharp, quantised
peaks. Both are exact consequences of admissibility, and both are what we measure. We do not touch the
asymptotic existence of long progressions; we map the finite-$X$ spectrum and verify its arithmetic.

\section{Setup: the AP constellation and its singular series}
A $k$-term arithmetic progression of primes is $p,\,p+\Delta,\,\dots,\,p+(k-1)\Delta$, all prime,
with common difference $\Delta\ge1$. We sieve all primes to $X=10^{8}$ ($5{,}761{,}455$ primes) and
count, for each $\Delta$,
\[
N_k(\Delta)=\#\bigl\{\,p\ge5\ \text{prime}\;:\;p+j\Delta\ \text{prime for }j=1,\dots,k-1,\ p+(k-1)\Delta\le X\,\bigr\}.
\]
The Hardy--Littlewood heuristic for the constellation $H_\Delta=\{0,\Delta,\dots,(k-1)\Delta\}$ gives
$N_k(\Delta)\sim \Sing_k(\Delta)\,c_k\,X/\log^k X$ for $\Delta\ll X$, with the $\Delta$-dependence
carried entirely by the singular series
\begin{equation}
\Sing_k(\Delta)=\prod_p \frac{1-\nu_p(\Delta)/p}{(1-1/p)^{k}},\qquad
\nu_p(\Delta)=\#\{\,H_\Delta \bmod p\,\}=
\begin{cases}1,& p\mid\Delta,\\[2pt]\min(k,p),& p\nmid\Delta.\end{cases}
\label{eq:Sk}
\end{equation}
The residue count $\nu_p$ is elementary: if $p\mid\Delta$ all terms collapse to one residue; if
$p\nmid\Delta$ the terms are $\Delta\cdot\{0,1,\dots,k-1\}\bmod p$, which are $\min(k,p)$ distinct
values. Because $N_k(\Delta)$ counts $\sim\!X$ starting points regardless of $\Delta$ (for
$\Delta\ll X$), the ratio of counts at two common differences is, to leading order, the ratio of their
singular series---which is what we measure.

\section{The selection rule: extinction off the primorial lattice}
Whenever $p\le k$ and $p\nmid\Delta$, the $k\ge p$ terms cover \emph{every} residue mod $p$, so
$\nu_p=p$ and the local factor $(1-\nu_p/p)=0$ annihilates $\Sing_k$. Hence:

\begin{proposition}[Admissibility selection rule]
$\Sing_k(\Delta)\ne0$ only if $\Delta$ is divisible by every prime $p\le k$, i.e.\
$\Delta\equiv0 \pmod{\mathrm{primorial}(k)}$, where $\mathrm{primorial}(k)=\prod_{p\le k}p$.
\end{proposition}

For $k=3$ and $k=4$ the lattice is $6\mathbb{Z}$; for $k=5$ it is $30\mathbb{Z}$. The arithmetic
reason is transparent: a term divisible by an unrequired small prime $p$ would have to equal $p$
itself to stay prime, which only the finitely many small progressions can do. Everything large is
extinguished. Figure~\ref{fig:fig1}(A) shows the measured $3$-AP spectrum as a comb: nonzero only on
multiples of $6$, dark between. The extinction is not approximate. Scanning \emph{every}
$\Delta\le1512$ for $k=3$, the number of off-lattice ($6\nmid\Delta$) common differences carrying a
single prime progression with $p\ge5$ is exactly zero. This is the lattice selection rule made
literal.

\section{The resonance enhancement spectrum}
On the admissible lattice the spectrum is far from flat. Adding a prime factor $q>k$ to $\Delta$
changes the local factor at $q$ from the generic $(1-k/q)/(1-1/q)^k$ (since $\nu_q=k$ when $q\ge k$
and $q\nmid\Delta$) to $(1-1/q)/(1-1/q)^k$ (since $\nu_q=1$ when $q\mid\Delta$). The ratio is the
universal enhancement
\begin{equation}
\frac{\Sing_k(\Delta\cdot q)}{\Sing_k(\Delta)}=\frac{1-1/q}{1-k/q}=\frac{q-1}{q-k},
\qquad q>k,\ q\nmid\Delta .
\label{eq:gain}
\end{equation}
Each extra prime $q$ dividing $\Delta$ multiplies the count by $(q-1)/(q-k)$: for $k=3$ a factor
$(q-1)/(q-3)$ (so $q=5$ doubles it, $q=7$ gives $\times1.5$); for $k=4$ a factor $(q-1)/(q-4)$ (so
$q=5$ \emph{quadruples} it); for $k=5$ a factor $(q-1)/(q-5)$ over primes $q\ge7$. The primorial of
$k$ is the price of admission; every prime above $k$ is a resonance.

Table~\ref{tab:res} and Figures~\ref{fig:fig1}(B), \ref{fig:fig2} put~\eqref{eq:gain} against the
data. The agreement is uniform to $\lesssim1\%$. At $k=3$: $\Delta=30$ (gains $5$) measures
$2.011$ versus $2.000$; $\Delta=210$ (gains $5,7$) measures $3.016$ versus $3.000$; $\Delta=2310$
(gains $5,7,11$) measures $3.768$ versus $3.750$. At $k=4$ the steeper law shows plainly: $\Delta=30$
measures $4.065$ versus $4.000$, $\Delta=2310$ measures $11.54$ versus $11.43$. At $k=5$, requiring
$30\mid\Delta$, $\Delta=210$ (gains $7$) measures $2.962$ versus $3.000$ and $\Delta=2310$ (gains
$7,11$) measures $4.955$ versus $5.000$. The enhancement quantises the spectrum into the discrete
bands visible in Figure~\ref{fig:fig1}(B), and the single-prime gains trace the curves
$(q-1)/(q-k)$ exactly (Fig.~\ref{fig:fig2}(B)).

\begin{table}[t]
\centering\small
\begin{tabular}{rrlrcc}
\toprule
$k$ & $\Delta$ & primes $q>k$, $q\mid\Delta$ & $N_k(\Delta)$ & measured ratio & predicted $\prod\frac{q-1}{q-k}$\\
\midrule
3 & 6    & ---     & 110392 & 1.0000 & 1.0000\\
3 & 30   & 5       & 222008 & 2.0111 & 2.0000\\
3 & 42   & 7       & 166258 & 1.5061 & 1.5000\\
3 & 66   & 11      & 138039 & 1.2504 & 1.2500\\
3 & 210  & 5,7     & 332970 & 3.0163 & 3.0000\\
3 & 330  & 5,11    & 277815 & 2.5166 & 2.5000\\
3 & 2310 & 5,7,11  & 415990 & 3.7683 & 3.7500\\
\midrule
4 & 6    & ---     & 9330   & 1.0000 & 1.0000\\
4 & 30   & 5       & 37927  & 4.0651 & 4.0000\\
4 & 210  & 5,7     & 75689  & 8.1124 & 8.0000\\
4 & 2310 & 5,7,11  & 107656 & 11.5387& 11.4286\\
\midrule
5 & 30   & ---     & 5523   & 1.0000 & 1.0000\\
5 & 210  & 7       & 16359  & 2.9620 & 3.0000\\
5 & 330  & 11      & 9060   & 1.6404 & 1.6667\\
5 & 2310 & 7,11    & 27367  & 4.9551 & 5.0000\\
\bottomrule
\end{tabular}
\caption{The resonance enhancement \eqref{eq:gain} against measured $k$-AP counts at $X=10^{8}$. The
ratio is relative to the primorial base ($\Delta=6$ for $k=3,4$; $\Delta=30$ for $k=5$). Agreement is
$\lesssim1\%$; the $k=5$ deviations are the largest because the base count is only $5{,}523$
(Poisson error $\approx1.3\%$).}
\label{tab:res}
\end{table}

\begin{figure}[t]
\centering
\includegraphics[width=\textwidth]{p30_fig1.pdf}
\caption{(A) The $3$-AP diffraction spectrum: $N_3(\Delta)$ versus $\Delta$ is a comb supported only on
multiples of $6$, extinguished elsewhere, with resonance peaks at $\Delta=30$ ($\times2$) and
$\Delta=210$ ($\times3$). (B) The admissible sub-spectrum $N_3(\Delta)/N_3(6)$ over all multiples of
$6$ (blue) against the relative Hardy--Littlewood singular series \eqref{eq:Sk} (red); the enhancement
quantises into discrete bands.}
\label{fig:fig1}
\end{figure}

\begin{figure}[t]
\centering
\includegraphics[width=\textwidth]{p30_fig2.pdf}
\caption{(A) The unified resonance law: measured $N_k(\Delta)/N_k(\mathrm{primorial})$ against the
predicted product $\prod_{q\mid\Delta,\,q>k}(q-1)/(q-k)$ for $k=3,4,5$, on the diagonal. (B) The
single-prime gain $(q-1)/(q-k)$ as a function of $q$ for each $k$ (curves), with the measured
single-prime points (markers) lying on them.}
\label{fig:fig2}
\end{figure}

\FloatBarrier
\section{Scope and honest ledger}
The experiment renders the count of prime arithmetic progressions, as a function of the common
difference, as a lattice spectrum with two exact features. The \emph{selection rule}---extinction off
the primorial sublattice $\mathrm{primorial}(k)\cdot\mathbb{Z}$---is admissibility, verified as a
literal zero for every off-lattice $\Delta\le1512$ at $k=3$. The \emph{resonance enhancement}---the
universal per-prime factor $(q-1)/(q-k)$ for primes $q>k$ dividing $\Delta$---is the relative
Hardy--Littlewood singular series, verified to $\lesssim1\%$ across $k=3,4,5$ and unified across $k$ by
a single closed form.

Neither feature is a new theorem; both follow from \eqref{eq:Sk}. What this paper contributes is the
spectral framing and the precision check at $X=10^{8}$, in the measurement spirit of the series: the
same singular-series machinery that gave the twin pair correlation of Part XXIX, here read along the
progression rather than across the pair. We are explicit about the boundary. The Green--Tao
theorem---that progressions of every length occur---concerns asymptotics no sieve can reach; at
$X=10^{8}$ the longest progressions are short, and our $k=5$ counts are already in the thousands. We
make no statement about the existence of long progressions, only about the arithmetic of the spectrum
where we can see it. Everything here equals the discrete sieve count, takes the Hardy--Littlewood
heuristic as input, and makes no claim about the infinitude of any constellation.

\paragraph{Data and code availability.} The prime sieve to $X=10^{8}$, the dense $3$-AP spectrum scan,
the cross-$k$ resonance table, and the figure scripts are openly available at
\url{https://github.com/Ruqing1963/6N-ap-resonance}.

\begin{thebibliography}{9}
\bibitem{HL} G.~H.~Hardy, J.~E.~Littlewood, \emph{Some problems of `Partitio Numerorum'; III}, Acta Math.\ \textbf{44} (1923), 1--70.
\bibitem{vdC} J.~G.~van der Corput, \emph{\"Uber Summen von Primzahlen und Primzahlquadraten}, Math.\ Ann.\ \textbf{116} (1939), 1--50.
\bibitem{GT} B.~Green, T.~Tao, \emph{The primes contain arbitrarily long arithmetic progressions}, Ann.\ of Math.\ \textbf{167} (2008), 481--547.
\bibitem{PI} R.~Chen, \emph{Conditional density on the 6N skeleton} (Part I), Zenodo, doi:10.5281/zenodo.20470367 (2026).
\bibitem{PXXIII} R.~Chen, \emph{Phase transition kinetics and topological event horizons at extreme $\omega$} (Part XXIII), Zenodo, doi:10.5281/zenodo.20586919 (2026).
\bibitem{PXXVI} R.~Chen, \emph{The second moment of $\omega$ under twin-centre conditioning on the 6N skeleton} (Part XXVI), Zenodo, doi:10.5281/zenodo.20593654 (2026).
\bibitem{PXXVIII} R.~Chen, \emph{Auditing the wind tunnel: the twin-enrichment closed form at reachable $\omega$} (Part XXVIII), Zenodo, doi:10.5281/zenodo.20596631 (2026).
\bibitem{PXXIX} R.~Chen, \emph{Twin counts in short intervals: singular-series rigidity and the density-gradient variance} (Part XXIX), Zenodo, doi:10.5281/zenodo.20597725 (2026).
\bibitem{review} R.~Chen, \emph{Arithmetic geodynamics on the 6N skeleton: a review of Parts I--XIX}, Zenodo, doi:10.5281/zenodo.20585301 (2026).
\end{thebibliography}

\end{document}
